% structure.tex — page geometry, running heads, part and chapter openings.
%
% Improvement #80. Everything here reads a token from tokens.tex; no length, size or grey
% is written twice.
%
% Before this item the book was A4 with a uniform 2.2 cm margin and fancyhdr's default
% heads, which means the gutter and the fore-edge were the same width and the text block
% drifted toward the spine on every recto. A bound book needs mirrored margins, and a
% reference book needs to say where you are on both sides of the spread.

% ---------------------------------------------------------------------------
% Geometry — mirrored, twoside
% ---------------------------------------------------------------------------
%
% `twoside` itself is set as a class option by build-book.sh. Setting it here would be too
% late: the class has already chosen its margin logic by the time this file is read.

\usepackage[
  \tokPaper,
  twoside,
  inner=\tokMarginInner,
  outer=\tokMarginOuter,
  top=\tokMarginTop,
  bottom=\tokMarginBottom,
  headheight=14pt,
  footskip=28pt,
]{geometry}

% A page that ends short is a page the reader thinks is the end of something. `\flushbottom`
% makes the block depth constant and pushes the slack into the inter-paragraph space, which
% is where \tokParSkip's elasticity is for.
\flushbottom

% A single line stranded at the top or bottom of a page is the most visible typesetting
% fault a reader can name without knowing the word for it. 10000 is TeX's "never".
\widowpenalty=10000
\clubpenalty=10000
\displaywidowpenalty=10000

\usepackage{needspace}
\usepackage{titlesec}
\usepackage{xstring}

% Pandoc's template sets `secnumdepth` to -\maxdimen when --number-sections is off, which
% is the right default for section headings and the wrong one for this book: it also
% unnumbers \part and \chapter, so the part opener printed no numeral and every chapter
% title lost its number. 0 numbers parts and chapters and stops there, which is exactly the
% book's own convention — "Chapter 14 — React Server Components", and no 14.2.3 anywhere.
\setcounter{secnumdepth}{0}

% ---------------------------------------------------------------------------
% Part openers
% ---------------------------------------------------------------------------
%
% The divider `collect-chapters.ts` emits reads `# Part 3 — The Modern Frontend Stack`,
% because that string has to work in EPUB and on the web as well, where nothing generates
% a part number. In print the numeral is set separately and large, so the prefix would be
% printed twice. `\StrBehind` cuts the title at the em dash and keeps the half that is a
% name; a divider with no dash — there are none today — falls back to the whole string
% rather than to nothing.

% `\StrBehind` assigns rather than expands, so it cannot run inside \markboth — TeX reads
% the `#1` in the mark's own definition and fails with "Illegal parameter number". The
% split therefore happens once, into \bookpartname, and everything else reads that.
% Only a divider that actually begins "Part " is split. The appendix's divider reads
% `# Appendix — DSA Patterns` and is not a numbered part of the book, so it keeps its whole
% title and prints no numeral — a big grey "X" above the DSA appendix would be claiming it
% is Part X, which BOOK-SPEC § 5 is explicit that it is not.
\newcommand{\bookpartname}{}
\newcommand{\bookpartnum}{}
\newcommand{\setpartname}[1]{%
  \IfBeginWith{#1}{Part }%
    {\StrBehind{#1}{— }[\bookparttail]%
     \IfStrEq{\bookparttail}{}%
       {\renewcommand{\bookpartname}{#1}}%
       {\let\bookpartname\bookparttail}%
     \renewcommand{\bookpartnum}{\thepart}}%
    {\renewcommand{\bookpartname}{#1}%
     \renewcommand{\bookpartnum}{}}%
}

\titleformat{\part}[display]
  {\raggedright\vspace*{6\baselineskip}}
  {\color{inklight}\sffamily\bfseries\fontsize{\tokPartNumberSize}{\tokPartNumberSize}\selectfont\bookpartnum}
  {14pt}
  {\color{ink}\sffamily\bfseries\fontsize{\tokPartTitleSize}{\tokPartTitleSize}\selectfont}
\titlespacing*{\part}{0pt}{0pt}{2\baselineskip}

% This has to come **after** \titleformat, which redefines \part itself — wrapping it
% before means titlesec overwrites the wrapper and nothing is stripped.
%
% Stripping at \part itself, rather than at each place the title is used, is what keeps the
% running head, the opener and the **table of contents** agreeing. The class hands the same
% argument to all three, so cleaning it once upstream is the only version that cannot drift.
\let\bookoldpart\part
\renewcommand{\part}[1]{\cleardoublepage\setpartname{#1}\bookoldpart{\bookpartname}}

% A part opener belongs on a recto — a reader turning pages sees it coming. `\cleardoublepage`
% is the book class's default for \part in twoside, and is restated in the redefinition
% above so that a future change to `openany` cannot silently move them.

% ---------------------------------------------------------------------------
% Chapter openings
% ---------------------------------------------------------------------------
%
% The source design spends 146pt before its first word. This spends 78pt:
%
%   eyebrow   8.5pt sans semibold caps, +140/1000 tracking, 55% K
%   12pt
%   number and title   22pt sans bold
%   14pt
%   rule   1.2pt
%   18pt
%
% The eyebrow is the part name, taken from `\leftmark`, so a reader who opens the book in
% the middle knows which of the nine parts they are in before reading a word of the
% chapter. The **deck** — the one-sentence promise under every chapter title — is still
% body text rather than part of this block: it is a blockquote in the markdown and pulling
% it into the opening needs a filter, which is #82's half of the work.

% Tracking comes from the font, not from a letterspacing package: `soul` cannot letterspace
% a macro like \leftmark without \mbox-ing it first, and microtype's \textls is pdftex-only.
% fontspec's LetterSpace applies the feature to the face itself, which is what \tokTracking
% means anyway — +140/1000 em.
\newcommand{\eyebrow}[1]{%
  {\labelfont\addfontfeature{LetterSpace=14.0}%
   \fontsize{\tokSmallSize}{\tokSmallSize}\selectfont\color{inkmid}%
   \MakeUppercase{#1}}%
}

\titleformat{\chapter}[display]
  {\raggedright}
  {\eyebrow{\leftmark}}
  {12pt}
  {\color{ink}\sffamily\bfseries\fontsize{\tokChapterSize}{\tokChapterSize}\selectfont
   \liningnums{\thechapter}\hspace{12pt}}
  [\vspace{14pt}{\color{ink}\titlerule[\tokAccentRule]}]
\titlespacing*{\chapter}{0pt}{0pt}{18pt}

% An unnumbered chapter — the preface, the glossary — has no part above it, so the eyebrow
% would print the previous part's name. It gets no eyebrow at all.
\titleformat{name=\chapter,numberless}[display]
  {\raggedright}
  {}
  {0pt}
  {\color{ink}\sffamily\bfseries\fontsize{\tokChapterSize}{\tokChapterSize}\selectfont}
  [\vspace{14pt}{\color{ink}\titlerule[\tokAccentRule]}]

% ---------------------------------------------------------------------------
% Section headings, and the guards that keep the closing blocks whole
% ---------------------------------------------------------------------------
%
% Every chapter ends on the same three headings — Key Takeaways, Interview Questions, What
% to Read Next. A heading that lands two lines above a page break puts its content on the
% next page and reads as a mistake. \needspace refuses to start one without four lines of
% room under it.

\titleformat{\section}
  {\needspace{4\baselineskip}\color{ink}\sffamily\bfseries\fontsize{13pt}{15pt}\selectfont}
  {}{0pt}{}
\titlespacing*{\section}{0pt}{1.4\baselineskip}{0.4\baselineskip}

\titleformat{\subsection}
  {\needspace{3\baselineskip}\color{ink}\sffamily\fontseries{sb}\selectfont\fontsize{11pt}{13pt}\selectfont}
  {}{0pt}{}
\titlespacing*{\subsection}{0pt}{1.1\baselineskip}{0.3\baselineskip}

% ---------------------------------------------------------------------------
% The table of contents
% ---------------------------------------------------------------------------
%
% 🔴 Improvement #77, found by measuring rather than by reading. `\@pnumwidth` is the box
% the page number sits in at the right edge of a contents line, and LaTeX sets it to
% 1.55em — about 15.5pt, sized for the three-digit folios of a normal book. This one runs
% past a thousand pages, and a four-digit lining number at 10pt is 20pt wide, so **every
% contents entry past page 999 hung its folio 4.5pt into the margin**: 499 of them, plus
% 206 more at 1.26pt for the three-digit entries. All of it under the 5pt the item's
% acceptance test allows, and all of it on the thirty pages a reader opens first.
%
% \@tocrmarg is the matching right margin for a title that wraps, and it has to stay
% ahead of \@pnumwidth or a wrapped line runs under the number instead of stopping
% before it.

% 🔴 Improvement #93, the other half of the same measurement. `\l@chapter` sizes the box
% the chapter *number* sits in — `\@tempdima` — at 1.5em, which is about 15pt and is
% sized for a book with two-digit chapters. This one numbers chapters continuously across
% nine parts and reaches three digits, so a "246" ran into the title beside it and TeX
% reported **412 overfull contents lines at 1.26pt each**. 2.4em clears a three-digit
% lining number at \tokBodySize with a word space after it.
%
% `\l@section` is left alone: the contents runs to depth 2, and a section number inherits
% its indent from the chapter box rather than setting its own.

\makeatletter
\renewcommand{\@pnumwidth}{2.2em}
\renewcommand{\@tocrmarg}{3.2em}

\renewcommand*\l@chapter[2]{%
  \ifnum \c@tocdepth >\m@ne
    \addpenalty{-\@highpenalty}%
    \vskip 1.0em \@plus\p@
    \setlength\@tempdima{2.4em}%
    \begingroup
      \parindent \z@ \rightskip \@pnumwidth
      \parfillskip -\@pnumwidth
      \leavevmode \bfseries
      \advance\leftskip\@tempdima
      \hskip -\leftskip
      #1\nobreak\hfil
      \nobreak\hb@xt@\@pnumwidth{\hss #2}\par
      \penalty\@highpenalty
    \endgroup
  \fi}
\makeatother

% ---------------------------------------------------------------------------
% Running heads and the folio
% ---------------------------------------------------------------------------
%
% Verso carries the part, recto the chapter, both at \tokMicroSize in 55% K caps. The
% folio sits on the outer edge so it stays outer whichever side the sheet is read from.
%
% `\leftmark` and `\rightmark` are what the book class fills with the part and chapter
% titles; the marks are reset here so that the part name survives into the chapter's
% eyebrow above.

\usepackage{fancyhdr}
\usepackage{xcolor}

% #1 has already been stripped by the \part redefinition above, so the mark only has to
% carry it. Splitting again here would be a second place for the rule to change.
\renewcommand{\partmark}[1]{\markboth{#1}{}}
\renewcommand{\chaptermark}[1]{\markright{#1}}

% #80 set the two marks above and stopped there, which left the recto head showing
% whichever `##` heading happened to end the page — "Common Mistakes" rather than the
% chapter title. The class's own \sectionmark still calls \markright, and it runs after
% \chaptermark does. Found by #81's specimen page; a section is not a location in this
% book, so it writes no mark at all.
\renewcommand{\sectionmark}[1]{}

\newcommand{\headfont}{\labelfont\fontsize{\tokMicroSize}{\tokMicroSize}\selectfont}

% ---------------------------------------------------------------------------
% Running page footer — copyright and the author's site on every page.
%
% Note this is a *running footer*, not chapter content: BOOK-SPEC.md non-negotiable
% #11 bans personal URLs from chapter bodies, and that still holds.
%
% EPUB gets the same line as `rights` metadata in scripts/book-meta.yaml. It is a
% reflowable format with no pages, so it has no footer to put this in.
% ---------------------------------------------------------------------------

\newcommand{\bookfooter}{%
  \headfont\color{inkmid}%
  \textcopyright{} Salman Rahman \quad\textperiodcentered\quad
  \href{https://www.salmanrahman.com/}{\textcolor{inkmid}{www.salmanrahman.com}}%
}

\pagestyle{fancy}
\fancyhf{}

\fancyhead[LE,RO]{\headfont\color{inkmid}\liningnums{\thepage}}
\fancyhead[RE]{\headfont\color{inkmid}\nouppercase{\leftmark}}
\fancyhead[LO]{\headfont\color{inkmid}\nouppercase{\rightmark}}
\fancyfoot[C]{\bookfooter}

\renewcommand{\headrulewidth}{\tokHairline}
\renewcommand{\footrulewidth}{\tokHairline}

% The book class forces \thispagestyle{plain} on the first page of every chapter and
% part. Without this, those pages would be the only ones with no copyright line.
\fancypagestyle{plain}{%
  \fancyhf{}%
  \fancyfoot[C]{\bookfooter}%
  \renewcommand{\headrulewidth}{0pt}%
  \renewcommand{\footrulewidth}{\tokHairline}%
}

% The title page uses the `titlepage` environment, which sets \thispagestyle{empty}.
% AtBeginEnvironment fires after that, so this wins. No rule here — a hairline under
% a title page reads as a mistake.
\usepackage{etoolbox}
\fancypagestyle{booktitle}{%
  \fancyhf{}%
  \fancyfoot[C]{\bookfooter}%
  \renewcommand{\headrulewidth}{0pt}%
  \renewcommand{\footrulewidth}{0pt}%
}
\AtBeginEnvironment{titlepage}{\thispagestyle{booktitle}}

% ---------------------------------------------------------------------------
% Genuinely blank versos
% ---------------------------------------------------------------------------
%
% The book class inserts a blank verso before every part and every chapter, and the stock
% \cleardoublepage leaves the current page style on it — so those pages came out carrying a
% folio, a running head and a copyright line under an otherwise empty leaf. That reads as a
% printing error rather than as a notice. This is the standard fix, and the only reason it
% needs \makeatletter is that \c@page and \if@twoside are internal names.

\makeatletter
\renewcommand{\cleardoublepage}{%
  \clearpage
  \if@twoside
    \ifodd\c@page\else
      \hbox{}\thispagestyle{empty}\newpage
    \fi
  \fi
}
\makeatother
